\documentclass[11pt]{article}
\usepackage{amsmath,amssymb,geometry}
\geometry{margin=25mm}
\begin{document}
\section*{Core model: air traffic flow management with enroute capacities (TFMP)}
Source: Bertsimas and Stock Patterson, Section 1.1, PDF pp.3--5
(printed pp.408--410), equations (2)--(8) and the TFMP objective.

\paragraph{Data.}
Let $F,K,J$ be flights, airports and enroute sectors, and let
$\mathcal T=\{1,\ldots,H\}$ be the finite time grid. Flight $f$ has the
fixed ordered path $P_f=(P(f,1),\ldots,P(f,N_f))$, with departure airport
$o_f=P(f,1)$ and destination airport $d_f^{loc}=P(f,N_f)$.
Each intermediate path element belongs to $J$. Let $\mathcal C$ contain
pairs $(f',f)$ flown consecutively by the same aircraft; their connecting
airport is $d_{f'}^{loc}=o_f$.
The scheduled departure/arrival are $d_f,r_f$, ground/air delay rates are
$c_f^g,c_f^a$, and $s_{f'}$ is the turnaround after $f'$.
The positive integer $l_{fj}$ is minimum traversal time from path element
$j$ to its successor. Capacities are $D_k(t)$ for departures, $A_k(t)$ for
arrivals, and $S_j(t)$ for simultaneous sector occupancy.
For each flight/path element, $T_f^j=\{\underline T_f^j,\ldots,\overline T_f^j\}$
is the supplied feasible arrival-time window.
\paragraph{Variables and boundary convention.}
$w_{ft}^j\in\{0,1\}$ indicates arrival at element $j$ by time $t$.
Use the source's window boundary values
\[
w_{ft}^j=0\quad(t<\underline T_f^j),\qquad
w_{ft}^j=1\quad(t\ge\overline T_f^j).
\]
Thus the last feasible time is fixed to one, ensuring every flight traverses
every element. Define the expressions
\[
D_f=\sum_{t\in T_f^{o_f}}t(w_{ft}^{o_f}-w_{f,t-1}^{o_f}),\quad
R_f=\sum_{t\in T_f^{d_f^{loc}}}t(w_{ft}^{d_f^{loc}}-w_{f,t-1}^{d_f^{loc}}),
\quad g_f=D_f-d_f,\quad a_f=R_f-r_f-g_f.
\]
\paragraph{One complete optimization model.}
\[
\min_w\ \sum_{f\in F}(c_f^g g_f+c_f^a a_f).
\]
Subject to
\begin{align}
\sum_{f:o_f=k}(w_{ft}^k-w_{f,t-1}^k)&\le D_k(t)
 &&k\in K,t\in\mathcal T,\tag{2}\\
\sum_{f:d_f^{loc}=k}(w_{ft}^k-w_{f,t-1}^k)&\le A_k(t)
 &&k\in K,t\in\mathcal T,\tag{3}\\
\sum_{\substack{f,i:P(f,i)=j\\j'=P(f,i+1),\ i<N_f}}
 (w_{ft}^{j}-w_{ft}^{j'})&\le S_j(t)
 &&j\in J,t\in\mathcal T,\tag{4}\\
w_{f,t+l_{fj}}^{j'}&\le w_{ft}^{j}
 &&f\in F,\ j=P(f,i),\ j'=P(f,i+1),\ i<N_f,\ t\in T_f^j,\tag{5}\\
w_{ft}^{k}&\le w_{f',t-s_{f'}}^{k}
 &&(f',f)\in\mathcal C,\ k=o_f,\ t\in T_f^k,\tag{6}\\
w_{ft}^{j}&\ge w_{f,t-1}^{j}
 &&f\in F,j\in P_f,t\in T_f^j,\tag{7}\\
w_{ft}^{j}&\in\{0,1\}
 &&f\in F,j\in P_f,t\in T_f^j.\tag{8}
\end{align}
All shifted indices use the boundary convention above. The objective is
algebraically the printed TFMP objective, including its schedule constants.
\end{document}
